% -*- coding: utf-8 -*-
% !TEX program = xelatex
\documentclass[plain,noanswer,medmath]{../../template/jnuexam}
\usepackage{../../template/kysx}

\begin{document}


\examtitle{1993}{2}


\loadproblems{1}{1993P1}%载入试卷一


\makepart{填空题}{本题共5小题，每小题3分，满分15分}%【同试卷一第一题】


\useproblem{1}{1}{1}


\useproblem{1}{1}{2}


\useproblem{1}{1}{3}


\useproblem{1}{1}{4}


\useproblem{1}{1}{5}


\makepart{选择题}{本题共5小题，每小题3分，满分15分}%【同试卷一第二题】


\useproblem{1}{2}{1}


\useproblem{1}{2}{2}


\useproblem{1}{2}{3}


\useproblem{1}{2}{4}


\useproblem{1}{2}{5}


\makepart{计算题}{本题共3小题，每小题5分，满分15分}%【同试卷一第三题】


\useproblem{1}{3}{1}


\useproblem{1}{3}{2}


\useproblem{1}{3}{3}


\makepart{计算题}{本题共3小题，每小题6分，满分18分}


\begin{problem}
设$z = x^3 f\left(xy,\frac{y}{x} \right)$，$f$具有连续二阶偏导数，
求$\frac{\pdz}{\pdy}$, $\frac{\pd^2 z}{\pd y^2}$
及$\frac{\pd^2 z}{\pdx\pd y}$．
\end{problem}


\useproblem{1}{4}{0}%【同试卷一第四题】


\begin{problem}
已知$R^3$的两个基为
$$\alpha_1 = \begin{pmatrix} 1 \\ 1 \\ 1 \end{pmatrix},\;
  \alpha_2 = \begin{pmatrix} 1 \\ 0 \\ -1 \end{pmatrix},\;
  \alpha_3 = \begin{pmatrix} 1 \\ 0 \\ 1 \end{pmatrix}
\quad \text{与}\quad
  \beta_1 = \begin{pmatrix} 1 \\ 2 \\ 1 \end{pmatrix},\;
  \beta_2 = \begin{pmatrix} 2 \\ 3 \\ 4 \end{pmatrix},\;
  \beta_3 = \begin{pmatrix} 3 \\ 4 \\ 3 \end{pmatrix},$$
求由基$\alpha_1,\;\alpha_2,\;\alpha_3$到基$\beta_1,\;\beta_2,\;\beta_3$的过渡矩阵．
\end{problem}


\makepart{}{本题满分7分}%【同试卷一第五题】

\useproblem{1}{5}{0}


\makepart{}{本题满分10分}%【同试卷一第六题】

\useproblem{1}{6}{0}


\makepart{}{本题满分8分}%【同试卷一第七题】

\useproblem{1}{7}{0}


\makepart{}{本题满分6分}%【同试卷一第八题】

\useproblem{1}{8}{0}


\makepart{}{本题满分6分}%【同试卷一第九题】

\useproblem{1}{9}{0}


\end{document}
